% titel: Blatt 3, A5: Ausbeute (Wasserstoff aus Zink)
% Teil II: der durchgerechnete Weg.

\subsection*{Gegeben und gesucht}
Reaktion: $\mathrm{Zn} + \mathrm{H_2SO_4} \rightarrow \mathrm{ZnSO_4} +
\mathrm{H_2}$ (1:1). Gegeben $m(\mathrm{Zn}) = \qty{80.0}{g}$, Ausbeute
$\eta = \qty{97.4}{\percent}$. Gesucht: tatsächlich gebildete
Masse $m(\mathrm{H_2})$.

Molare Massen (gerundet): $M(\mathrm{Zn}) \approx \qty{65.38}{g/mol}$,
$M(\mathrm{H_2}) \approx \qty{2.016}{g/mol}$.

\subsection*{Ansatz}
Stoffmenge des Zinks, wegen 1:1 gleiche theoretische Wasserstoff-Stoffmenge,
daraus theoretische Masse; mit $\eta$ multiplizieren.

\subsection*{Rechnung}
\begin{align*}
  n(\mathrm{Zn}) &= \frac{80.0}{65.38} = 1.224\ \mathrm{mol}
  \;\Rightarrow\; n_{\text{theor}}(\mathrm{H_2}) = 1.224\ \mathrm{mol}, \\[2pt]
  m_{\text{theor}}(\mathrm{H_2}) &= 1.224 \times 2.016 = 2.467\ \mathrm{g}, \\
  m_{\text{real}}(\mathrm{H_2}) &= 0.974 \times 2.467
  \approx \qty{2.40}{g}.
\end{align*}

\begin{ergebnis}
Aus $\qty{80.0}{g}$ Zink entstehen bei einer Ausbeute von
$\qty{97.4}{\percent}$ etwa $m(\mathrm{H_2}) \approx \qty{2.40}{g}$
Wasserstoff.
\end{ergebnis}

\begin{kontrolle}
Bei 100\% wären es $2.47\ \mathrm{g}$; $97.4\%$ davon sind $2.40\ \mathrm{g}$,
also minimal weniger als die theoretische Menge --- plausibel. Querprobe der
Grössenordnung: $1.224\ \mathrm{mol}$ eines $\mathrm{H_2}$ (Molmasse 2 g/mol)
ergeben rund $2.4\ \mathrm{g}$.
\end{kontrolle}

\begin{merksatz}
Erst die 100\%-Menge berechnen (Eduktmenge $\rightarrow$ Stöchiometrie
$\rightarrow$ Produktmenge), dann mit der Ausbeute als Bruch multiplizieren.
\end{merksatz}
